\def \Subject {سوال دوم - حریم خصوصی در یادگیری ماشین و یادگیری فدرال  }
\def \Session { دو}
% \setcounter{chapter}{\Session-1}
\renewcommand{\chaptername}{سوال}
\chapter{\Subject}

\section{بخش اول}

\subsection{ زیربخش اول - حساسیت و مکانیزم لاپلاس}
\subsubsection{الف}
\begin{enumerate}
	\item 
	\textbf{میانگین مصرف:}
	در اینجا بیشترین تفاوت وقتی است که یک خانوار بدون مصرف با یک خانوار که بیشترین مصرف را دارد جایگزین شود. در اینصورت داریم:
	\begin{latin}
		~\\
		$
		max(q_1(D) - q_1(D')) = \cfrac{sum + 100}{1000} - \cfrac{sum}{1000} = 0.1
		$
	\end{latin}
		\item 
	\textbf{کل مصرف}
این مشابه بخش میانگین است با این تفاوت که نیازی نیست به تعداد تقسیم کنیم. پس داریم:
	\begin{latin}
	~\\
	$
	max(q_2(D) - q_2(D')) = sum + 100 - sum  = 100
	$
\end{latin}
\item 
\textbf{تعداد خانوار در هر الگوی مصرف:}
در اینجا بیشترین تفاوت وقتی است که الگوی مصرف دو خانواری که باهم تفاوت دارند، هیچ اشتراکی نداشته باشد. پس یعنی بردار تعداد الگوی مصرف در ۴ الگو، اختلاف یک دارد (یا در دو الگو اختلاف ۲) پس یعنی مقدار حساسیت برابر با ۴ است. 
	
\end{enumerate}
\subsubsection{ب}

	برای مکانیزم لاپلاس داشتیم:
	\begin{latin}
		~\\
		$
		b = \cfrac{\Delta}{\epsilon} 
		\xrightarrow{\epsilon = 0.2} b = 5\Delta
		$
	\end{latin}
	با جایگذاری 
	$\Delta$
	برای هر مورد داریم:
	\begin{enumerate}
	\item 
	\textbf{میانگین مصرف:}
	\begin{latin}
		~\\
		$
		b = 5\Delta = 0.5
		$
	\end{latin}
	\item 
	\textbf{کل مصرف}

	\begin{latin}
		~\\
		$
b = 5\Delta = 500
		$
	\end{latin}
	\item 
	\textbf{تعداد خانوار در هر الگوی مصرف:}
\begin{latin}
	~\\
	$
	b = 5\Delta = 20
	$
\end{latin}
	
\end{enumerate}
\subsection{زیربخش دوم - انتشار خصوصی داده‌ها}
\subsubsection{الف}
ازآنجایی که نویز اعمال شده روی هر کوئری مستقل از سایر کوئری‌ها می‌باشد، این ناسازگاری می‌تواند رخ دهد. در مثال سوال، چون نویز میانگین مستقل از نویز جمع مصرف است، مقدار میانگین گزارش داده شده با جمع گزارش داده شده ناسازگار است.

در اینجا چون می‌دانیم که جمع مصرف هزار برابر میانگین مصرف است، می‌توان این قید را اجبار کرد که پس از اعمال نویز بتوان این قید را ارضا کرد. 
\subsubsection{ب}
\begin{latin}
	~\\
	$
	b = 2.5 \Delta = 10 
	$
\end{latin}
\subsubsection{ج}
در این حالت بجای اینکه 
$\Delta$
برابر با ۴ شود،‌برابر با یک می‌شود. پس داریم:
\begin{latin}
	~\\
	$
	b = 2.5 \Delta = 2.5
	$
\end{latin}
حال اگر نویزهای داده‌شده را روی داده اعمال کنیم داریم:
\begin{latin}
	~\\
	$
	v = \begin{bmatrix}
		310 & 295 & 180 & 160 & 90 & 75
	\end{bmatrix}
	\\
	noise = \begin{bmatrix}
		-4 & 18 & 9 & -6 & 21 & 2
	\end{bmatrix}
	\\
	\Rightarrow 
	v + noise = \begin{bmatrix}
		306 & 313 & 189 & 154 & 111 & 77
	\end{bmatrix}
	\Rightarrow argmax(v + noise) = 2
	$
\end{latin}
یعنی در داده‌های اصلی،‌داده با بیشترین تکرار الگوی مصرف یک بود ولی در داده‌های نویزی الگوی مصرف دو به عنوان بیشینه معلوم شد و خروجی داده می‌شود. 

این موضوع به این بده بستان اشاره دارد که هرچقدر بخواهیم بیشتر حریم خصوصی را رعایت کنیم، داده‌ها بیشتر نویزی می‌شوند و در نتیجه اطلاعات و الگوهای داده‌ها بیشتر خراب می‌شود پس دقت مدل (و یا عملکرد سیستم استفاده کننده بصورت کلی) افت می‌کند.

\subsection{زیربخش سوم - تخصیص بودجه حریم خصوصی و ترکیب مکانیزم‌ها}
\subsubsection{الف}
از آنجایی که دقت میانگین مصرف مهم‌تر است پس یعنی بالاجبار حریم خصوصی در آن بیشتر فدا می‌شود پس سهم بیشتری از 
$\epsilon$
باید به میانگین مصرف داده شود. تقسیم‌بندی پیشنهادی من در جدول 

\begin{table}[H]
	\centering
	\caption{تقسیم $\epsilon$ بین مکانیزم‌ها}
	\label{tab:eps_tokhs}
	

		\begin{tabular}{| r | c|}
			\hline
			\textbf{مکانیزم} & $\epsilon$ \\
			\hline
			میانگین مصرف & $0.4$ \\
			پرتکرار ترین الگوی مصرف & $0.1$ \\
			الگوی مصرف شماره ۶ & $0.1$ \\
			\hline
		\end{tabular}

	
\end{table}

\subsubsection{ب}
داریم:
\begin{latin}
	~\\
	$
	b = \cfrac{\Delta}{\epsilon} = \cfrac{0.1}{0.4} = 0.25
	$
\end{latin}

\subsubsection{ج}
داریم:
\begin{latin}
	~\\
	$
	l_2 = max(||f(D) - f(D')||_2) = \sqrt{1} = 1  \qquad \epsilon = 0.1 \qquad \delta = 0.1
	\\
	\sigma \geq \cfrac{\sqrt{2ln (\cfrac{1.25}{\delta})}}{\epsilon}  
	\xrightarrow{c \approx 2.25} 
	\sigma \geq \cfrac{2.25}{0.1} 
	\Rightarrow \sigma \geq 22.5
	$
\end{latin}
\subsubsection{د}
طبق قانون 
\lr{composition}
مقدار 
$\epsilon$
ها باهم جمع می‌شود. پس یعنی مقدار 
\lr{privacy loss}
برابر با 
$7 * 0.6  = 4.2$
می‌شود. اگر بخواهیم 
\lr{privacy loss}
کل هفته برابر با 
$0.6$
باقی بماند، باید 
\lr{privacy loss}
مجموع
هر روز را برابر با 
$\cfrac{0.6}{7}$
درنظر بگیریم. این موضوع باعث می‌شود تا مقدار نویز خیلی زیاد شود و 
\lr{utility }
کاهش می‌یابد.

\section{بخش دوم - سوالات مفهومی}
\subsection{زیربخش اول - حریم خصوصی، استحکام و یادگیری فدرال}
\subsubsection{الف}
خیر نمی‌توان. برای مثال در درس حمله‌های مبتنی بر گرادیان را داشتیم که تقریبا ورودی اولیه را بصورت کامل بازسازی می‌کردند. 

بهبود عملی حریم خصوصی به معنی الگوریتم‌هایی است که در عمل حریم خصوصی را بهبود می‌دهند ولی اثبات ریاضی برای حد و حدود آن ندارند. در مقابل تضمین رسمی حریم خصوصی به معنی الگوریتم و یا اثبات‌هایی است که کران حریم خصوصی را بطور ریاضی نشان و ثابت می‌کنند (این بخش رو من از تو مقاله پیدا نکردم. با استفاده از ai خوندم که هر کدوم از الگوریتم ها تو کدوم دسته است.)

\subsubsection{ب}
حریم خصوصی به این معنا است که مدل از اطلاعات ورودی خود محافظت کند و مهاجم نتواند راجع به داده آموزشی اطلاعاتی کسب کند. ولی مقاوم بودن به این معنا است که مدل نسبت به حملات مخرب، مثلا حملات 
\lr{backdoor}
و یا 
خصمانه مقاوم باشد.

مثلا با
\lr{LDP}
اگر کلاینتی 
\lr{LDP}
را اعمال نکند فضای حل بزرگ‌تری خواهد داشت و می‌تواند تاثیر بیشتری بر مدل بگذارد. پس اگر یک مهاجم که قصد حمله 
\lr{backdoor}
داشته باشد داشته باشیم، می‌تواند راحت تر مدل را آلوده کند. در اینجا حریم خصوصی همچنان حفظ شده ولی این حفظ حریم خصوصی باعث شده است تا مدل کمتر مقاوم باشد.

\subsubsection{ج}
خیر. زیرا در درس داشتیم که با صعود گرادیان، و معکوس کردن فرآیند 
\lr{Backpropagation}
می‌توان ورودی را تقریبا کامل بازسازی کرد. هرچند که مدل 
\lr{B}
باز داده کمتری می‌دهد و همچنین راه را برای حفظ حریم خصوصی بیشتر باز می‌گذارد، ولی فقط همینکه داده نفرستیم و به روزرسانی وزن‌ها فرستاده شود، تقریبا هیچ فایده‌ای در حریم خصوصی نخواهد داشت.
\subsection{زیربخش دوم - حملات حریم خصوصی و مسموم سازی در یادگیری فدرال}
\subsubsection{الف}
ناقص 
\subsubsection{ب}
همانطور که در مقاله نیز اشاره می‌شود، هرکاری که با مسموم کردن داده قابل انجام است با مسموم کردن مدل نیز قابل انجام است. ولی مسموم کردن مدل انعطاف بیشتری دارد. مثلا برای مسموم کردن داده اولا نیاز به داده‌های خراب زیادی است تا بتوان علیرغم سایر کلاینت‌ها همچنان روی مدل تاثیر گذاشت، دوما سرور مرکزی ممکن است با کارهایی مثل 
\lr{byzantine tolerance}
داده‌های خراب را تشخیص دهد و سوما چون حمله مبتنی بر داده است، تعداد و یا توان مهاجمان باید زیاد باشد. چهارما اگر مکانیزمی مثل 
\lr{LDP}
داشته باشیم، در مسموم‌سازی داده اثر آن می‌تواند کم رنگ شود ولی در مسموم‌سازی مدل بسادگی می‌توان نویز را نادیده گرفت.
ولی در مسموم کردن مدل، چون می‌توان مستقیما خود مقدار مقادیر به روز رسانی را تغییر داد، می‌توان با اعمال یک 
\lr{norm}
از پرت بودن مقادیر نسبت به سایر کلاینت‌ها جلوگیری کرد، یا با زیاد کردن ضریب وزن‌ها تاثیر مهاجم را بیشتر کرد و نیازی هم به داده مجزا ندارد.

\subsubsection{ج}
\begin{itemize}
	\item 
	ایده کلی حمله مشابه حمله صعود گرادیان است. منتها در حالت کلی، از آنجایی که ممکن است اثر وجود یک عضو در مجموعه داده کلاینت، توسط به روزرسانی‌های سایر کلاینت‌ها کم شود، سرور مرکزی به یک کلاینت خاص (که کلاینت a می‌نامیم) وزن‌های خود کلاینت a را می‌دهد. در  نتیجه مدل a فقط روی داده‌های خودش آموزش می‌بیند و تاثیر حضور و یا عدم حضور یک داده در مجموعه داده آن بیشتر است. 
	
	در قدم بعدی وقتی سرور مرکزی می‌خواهد ببیند داده x در مجموعه داده کلاینت a بوده یا نه، x را به مدل ورودی می‌دهد و گرادیان را حساب می‌کند و وزن‌ها را در جهت افزایش آن حرکت می‌دهد. اگر مدل قبلا این ورودی را دیده باشد، تغییر وزن‌ها معنی‌دار می‌شود ولی اگر این داده در مجموعه داده آموزشی نباشد، مقدار تغییر وزن‌ها کوچک است.
	\item 
	اگر به شمای گرافیکی تابع یادگرفته شده توسط شبکه‌های عصبی نگاه کنیم، عمدتا به این شکل هست که یک رویه در مکان‌هایی مقدار خروجی پیک زده (به نوعی شبیه ترکیب چند توزیع نرمال) این کار به اینصورت است که وزن‌های اولیه شبکه بسیار کوچک هستند و می‌توان شبکه را به یک لایه نسبتا هموار تشبیه کرد. وقتی برای یک ورودی وزن‌ها میخواهند به روز رسانی شوند، این رویه دچار قله/دره می‌شود. حال اگر نمونه‌ای که به مدل ورودی داده می‌شود نمونه جدیدی باشد، انگار یک نقطه روی بخش هموار رویه را ورودی داده‌ایم و مقدار آن کم است (شیب در آن ناحیه خیلی زیاد نیست.) ولی وقتی یک داده که از قبل دیده شده ورودی داده شود، مدل با تغییر وزن می‌خواهد از یکی از این قله/دره‌ها در بیاید پس تغییر زیاد است.
	 
\end{itemize}
\subsection{زیربخش سوم}
\subsubsection{خطوط پایه}
\begin{itemize}
	\item 
	\textbf{\lr{Norm Bounding:}}
	این روش در هر مرحله به روزرسانی کلاینت‌ها را دریافت می‌کند و درصورتی که نرم یک به روزرسانی از یک حد آستانه بیشتر باشد یا آن را اصلا اعمال نمی‌کنیم و یا مقادیر آن را کاهش می‌دهیم تا نرم آن در حد آستانه قرار بگیرد. 
	\item \textbf{\lr{WeakDP:}}
	این روش ابتدا 
	\lr{Norm Bounding}
	را انجام می‌دهد و سپس یک نویز گاوسی با واریانس ثابت را روی نتیجه اعمال می‌کند. ولی از آنجایی که واریانس از روی 
	$\epsilon$
	بدست نمی‌آید، مقدار 
	\lr{privacy loss}
	در این مکانیزم خیلی زیاد است و عملا 
	\lr{differential privacy }
	خاصی ارائه نمی‌دهد. 
\end{itemize}
\subsubsection{\lr{CDP} و \lr{LDP}}
\begin{itemize}
	\item 
	\textbf{\lr{CDP:}}
	در این روش مکانیزم اضافه کردن نویز برای 
	\lr{Differential Privacy}
	در سرور مرکزی و روی همه داده‌ها اعمال می‌شود. 
	\item 
	\textbf{\lr{LDP:}}
	در این روش مکانیزم اضافه کردن نویز برای 
	\lr{Differential Privacy}
	در کلاینت‌ها و قبل از ارسال وزن‌ها به سرور مرکزی انجام می‌شود.
\end{itemize}

تفاوت اصلی در این است که 
\lr{LDP}
برای دفاع در برابر حمله از سمت سرور مرکزی و 
\lr{CDP}
برای دفاع در برابر کلاینت‌های مخرب است.  مثلا حمله ویژگی، که از سمت یک کلاینت مخرب است با 
\lr{LDP}
فایده‌ای ندارد ولی 
\lr{CDP}
بصورت تئوری می‌تواند در برابر این حمله مقداری دفاع کند. (هرچند که نتیجه خوبی نمی‌دهد.)

برای حمله‌هایی که نه از سمت کلاینت و نه از سمت سرور هستند هر دو رویکرد نتیجه خوبی دارند. 

\subsubsection{ناکارآمدی \lr{CDP}}
در این حملات 
\lr{CDP }
توانایی دفاع را دارد. ولی چون مقدار
$\epsilon$
کم است، باعث می‌شود تا نویز شدید شود و 
\lr{utility}
خیلی کاهش یابد. عملا دلیل ضعف مدل با 
\lr{CDP}
و ضعیف بودن 
\lr{weak DP}
ریشه یکسانی دارد. تقابل 
\lr{privacy}
و
\lr{utility}

\subsubsection{طراحی سیستم پزشکی}
اگر چالشی برای کلاینت مخرب نداریم (که با توجه به اینکه ده بیمارستان بزرگ کلاینت ها هستند احتمالا نداشته باشیم) استفاده از 
\lr{CDP}
بهتر است. 
\lr{CDP}
هم 
\lr{utility}
بهتری دارد و هم در برابر حملات، نرخ حمله موفق کمتری دارد.

\section{بخش سوم - سوالات پیاده‌سازی}











































\subsection{ زیربخش اول: آماده‌سازی داده‌ها}

\subsubsection{بارگذاری و بررسی اولیه داده}

دیتاست FE
\begin{table}[H]
	\centering
	\caption{آمار دیتاست \lr{FEMNIST}}
	\label{tab:femnist_features}
	\begin{tabular}{|l|r|}
		\hline
		\textbf{ویژگی} & \textbf{مقدار} \\
		\hline
		تعداد کل نمونه‌ها & ۸۱۴,۲۷۷ \\
		تعداد کلاس‌ها  & ۶۲ \\
		تعداد نویسندگان  & ۳,۵۹۷ \\
		تعداد نویسندگان عضو  & ۳,۳۹۷ \\
		تعداد نویسندگان غیرعضو  & ۲۰۰ \\
		تعداد نمونه‌های آموزشی  & ۷۶۹,۶۳۴ \\
		تعداد نمونه‌های تست  & ۴۴,۶۴۳ \\
		\hline
	\end{tabular}
\end{table}

همان‌طور که در جدول مشخص است، دیتاست شامل تنوع بالایی از نویسندگان و کلاس‌ها است. تصاویر نمونه‌ای از این دیتاست در شکل \ref{fig:sample_images} نمایش داده شده است که نشان می‌دهد تصاویر شامل اعداد و حروف دست‌نوشته‌ی مختلف با اندازه‌ی ۲۸ در ۲۸ پیکسل هستند. توزیع تعداد نمونه‌ها در بین نویسندگان نیز در شکل \ref{fig:writer_distribution} نشان داده شده است که حاکی از توزیع نابرابر است و برخی نویسندگان تعداد نمونه‌های بسیار بیشتری نسبت به دیگران دارند. این با موضوع مقاله که تفاوت حمله‌های 
\lr{record level}
و
\lr{subject level}
در 
\lr{federated learning }
است مرتبط است.

\begin{figure}[H]
	\centering
	\includegraphics[width=0.8\textwidth]{images/sample_images.png}
	\caption{نمونه تصاویر دیتاست FEMNIST به همراه برچسب واقعی}
	\label{fig:sample_images}
\end{figure}

\begin{figure}[H]
	\centering
	\includegraphics[width=0.8\textwidth]{images/writer_distribution.png}
	\caption{توزیع تعداد نمونه‌ها به ازای هر نویسنده}
	\label{fig:writer_distribution}
\end{figure}

\subsubsection{تعریف مجموعه‌های Member و Non-Member}

دیتاست به صورت تصادفی به دو مجموعه تقسیم شدد. ۳,۳۹۷ نویسنده به عنوان Member Subjects و ۲۰۰ نویسنده به عنوان Non-Member Subjects در نظر گرفته شدند.  جدول \ref{tab:member_nonmember} تعداد نویسندگان و تصاویر مربوط به هر مجموعه را نشان می‌دهد.

\begin{table}[H]
	\centering
	\caption{تعداد نویسندگان و تصاویر در مجموعه‌های Member و Non-Member}
	\label{tab:member_nonmember}
	\begin{tabular}{|l|r|r|}
		\hline
		\textbf{مجموعه} & \textbf{تعداد نویسندگان} & \textbf{تعداد تصاویر} \\
		\hline
		Member Subjects & ۳,۳۹۷ & ۷۶۹,۶۳۴ \\
		Non-Member Subjects & ۲۰۰ & ۴۴,۶۴۳ \\
		\hline
	\end{tabular}
\end{table}

\subsubsection{ایجاد ساختار Cross-Silo}

داده‌ها به  ۸ گروه مستقل تقسیم شدند. روش توزیع به این گونه انجام شد که برای هر نویسنده از روی یک توزیع شبیه به توزیع نمایی (ولی بصورت گسسته و clip شده بین ۱ تا ۶) معلوم شد که باید داده‌های این نویسنده بین چند سازمان پخش شود و سپس این پخش شدن انجام شد. جدول \ref{tab:silo_stats} و شکل \ref{fig:silo_share} آمار هر سازمان را نشان می‌دهند.

\begin{table}[H]
	\centering
	\caption{توزیع داده‌ها بین سازمان‌ها}
	\label{tab:silo_stats}
	\begin{tabular}{|l|r|r|r|}
		\hline
		\textbf{سازمان} & \textbf{تعداد تصاویر} & \textbf{درصد از کل داده‌ها} & \textbf{تعداد نویسنده‌ها} \\
		\hline
		Silo 0 & ۱۹۱,۰۴۷ & ۱۲.۴\% & ۳,۳۹۷ \\
		Silo 1 & ۱۹۱,۱۰۱ & ۱۲.۴\% & ۳,۳۹۷ \\
		Silo 2 & ۱۹۱,۳۱۰ & ۱۲.۴\% & ۳,۳۹۷ \\
		Silo 3 & ۱۹۱,۲۶۶ & ۱۲.۴\% & ۳,۳۹۷ \\
		Silo 4 & ۱۹۱,۷۱۴ & ۱۲.۵\% & ۳,۳۹۷ \\
		Silo 5 & ۱۹۱,۴۶۲ & ۱۲.۴\% & ۳,۳۹۷ \\
		Silo 6 & ۱۹۱,۲۷۷ & ۱۲.۴\% & ۳,۳۹۷ \\
		Silo 7 & ۱۹۱,۲۵۸ & ۱۲.۴\% & ۳,۳۹۷ \\
		\hline
	\end{tabular}
\end{table}

\begin{figure}[H]
	\centering
	\includegraphics[width=0.6\textwidth]{images/share_of_each_silo.png}
	\caption{سهم هر سازمان از کل داده‌های آموزشی}
	\label{fig:silo_share}
\end{figure}

\subsubsection{تحلیل آماری و مصورسازی}

تحلیل آماری انجام شده بر روی ساختار داده‌ها، نشان‌دهنده‌ی ویژگی‌های مهمی است که در شکل‌های \ref{fig:writer_silo_count} و \ref{fig:multi_silo} قابل مشاهده است.

\begin{figure}[H]
	\centering
	\includegraphics[width=0.6\textwidth]{images/writer_silo_count_distribution.png}
	\caption{توزیع تعداد سازمان‌هایی که هر Subject در آن حضور دارد}
	\label{fig:writer_silo_count}
\end{figure}

\begin{figure}[H]
	\centering
	\includegraphics[width=0.6\textwidth]{images/single_vs_multi_silo.png}
	\caption{درصد Subject‌های حاضر در یک سازمان در مقابل چند سازمان}
	\label{fig:multi_silo}
\end{figure}

همان‌طور که در شکل \ref{fig:writer_silo_count} مشاهده می‌شود، تعداد سازمان‌ها به ازای هر نویسنده شبیه به یک توزیع نمایی است. شکل \ref{fig:multi_silo} نشان می‌دهد که حدود ۵۲\% از نویسنده‌ها در یک سازمان و ۴۸\% در چند سازمان توزیع شده‌اند.


\subsection{زیربخش دوم: آموزش مدل }

. یک مدل CNN ساده با معماری زیر برای آموزش طراحی گردید:

\begin{itemize}
	\item دو لایه کانولوشن با ۱۶ و ۳۲ فیلتر به ترتیب و kernel size برابر با ۳
	\item لایه‌های ReLU و MaxPooling برای هر لایه کانولوشن
	\item لایه AdaptiveAvgPool2d برای کاهش ابعاد
	\item دو لایه Fully Connected با ۱۲۸ و ۶۲ نرون (به تعداد کلاس‌ها)
\end{itemize}

پارامترهای آموزش به صورت زیر تنظیم شدند:
\begin{itemize}
	\item تعداد دورهای ارتباطی: ۵
	\item تعداد Epochهای محلی: ۱
	\item نرخ یادگیری: ۰.۰۵
	\item تابع هزینه: CrossEntropyLoss
	\item بهینه‌ساز: SGD با momentum 0.9
\end{itemize}

فقط از داده‌های 
\lr{member}
برای آموزش استفاده شد و از داده‌های 
\lr{non member}
برای آموزش استفاده‌ای نشد. البته برای 
\lr{evaluation}
عملکرد مدل روی این داده‌ها حساب شد ولی مدل آموزش داده نشد.

مکانیزم تعامل بین سازمان‌ها و Server به این صورت بود که در هر دور، مدل جهانی به هر سازمان ارسال شده و سازمان‌ها مدل را بر روی داده‌های محلی خود به تعداد Epochهای مشخص آموزش می‌دهند. سپس وزن‌های به‌روزرسانی شده به سرور ارسال شده و سرور با میانگین‌گیری از وزن‌های تمام سازمان‌ها، مدل جهانی جدید را تولید می‌کند.

نتایج آموزش در شکل \ref{fig:fed_learning_metrics} نشان داده شده است. دقت مدل بر روی داده‌های آموزشی Member Subjects از 

$0.0905$
در دور اول به
$0.8004$
 در دور پنجم رسید که نشان‌دهنده‌ی همگرایی مناسب مدل است. همچنین دقت مدل بر روی داده‌های تست Non-Member Subjects که در فرآیند آموزش شرکت نداشتند، از
$0.0930$
  به
$0.8042$
   افزایش یافت.

\begin{figure}[H]
	\centering
	\includegraphics[width=0.9\textwidth]{images/federated_learning_metrics.png}
	\caption{دقت و تابع هزینه مدل در طول دورهای آموزش فدرال}
	\label{fig:fed_learning_metrics}
\end{figure}




\begin{table}[H]
	\centering
	\caption{اظهارنامه استفاده از هوش مصنوعی}
	\label{tab:ai_declaration}
	\begin{tabular}{|p{3cm}|p{4cm}|p{8cm}|}
		\hline
		\textbf{ابزار} & \textbf{برای چه کاری استفاده شد} & \textbf{شرح استفاده} \\
		\hline
		ChatGPT & ریویو و درک مقاله‌ها & برای مرور سریع مقاله‌های مرتبط با COMPAS، Membership Inference، و Differential Privacy از AI کمک گرفتم تا مفاهیم کلیدی را سریع‌تر درک کنم و به ساختار مقاله مسلط شوم. \\
		\hline
		ChatGPT & پیاده‌سازی فیلترهای ProPublica & از AI در مورد فیلترهای استاندارد ProPublica برای پالایش دیتاست COMPAS سوال پرسیدم و منطق فیلتر کردن داده‌ها را بررسی کردم. \\
		\hline
		ChatGPT & تبدیل جدول‌ها به LaTeX & برای تبدیل جدول‌های خروجی pandas به کد LaTeX و فرمت‌دهی مناسب جداول در گزارش از AI کمک گرفتم. \\
		\hline
		ChatGPT & بهبود ظاهر جداول و نمودارها & برای زیباسازی جداول (افزودن خطوط، تنظیم ستون‌ها) و تنظیم اندازه و کیفیت نمودارها (افزودن برچسب، عنوان، ذخیره با dpi مناسب) از AI کمک گرفتم. \\
		\hline
		ChatGPT & ریویو و ویرایش نگارشی گزارش & پس از نوشتن پیش‌نویس گزارش، از AI برای ویرایش نگارشی، بهبود روان‌نویسی و یکدست‌سازی لحن متن استفاده کردم. \\
		\hline
		ChatGPT & کد ریویو & کدهایی که من برای آموزش مدل زدم یه مقدار زیادی پیچیده بودن. برای ریویوش از gpt کمک گرفتم\\
		\hline
	\end{tabular}
\end{table}

